\Chapter{REVUE DE LITTÉRATURE}\label{sec:RevLitt}

Cette revue de littérature présente une analyse structurée de l'état de l'art pertinent pour la génération de trajectoires robotisées sûres en environnement dynamique et non structuré. Elle est organisée selon la logique du pipeline « perception, planification, action » proposé dans ce travail : les méthodes de génération de mouvement, dont le \textit{Flow Matching} retenu comme planificateur nominal ; la perception visuelle et la mémoire spatiale qui fournissent le conditionnement géométrique ; les représentations géométriques permettant d'évaluer les collisions ; les lois de commande sécuritaires réactives ; et enfin les stratégies de couplage entre modèles génératifs et filtres de sécurité. Chaque section expose les contributions existantes, en dégage les forces et les limites, et prépare le positionnement de l'architecture découplée développée au chapitre suivant.

% =================================================================
\section{Génération de trajectoires robotisées}
% =================================================================

\noindent\textbf{Objectif de la section :} Analyser le développement théorique et les limites connues des méthodes de génération de mouvements afin de situer le choix du \textit{Flow Matching} conditionnel comme planificateur d'intention nominale.

\subsection{Planification classique par échantillonnage et optimisation}
La génération de mouvements sans collision a d'abord été abordée par des algorithmes déterministes qui raisonnent directement dans l'espace des configurations, sans apprentissage préalable. Deux grandes familles se dégagent, selon qu'elles explorent cet espace par tirages aléatoires ou qu'elles raffinent une trajectoire par optimisation continue.

\paragraph{Méthodes par échantillonnage.} L'algorithme \textit{Rapidly-exploring Random Tree star} (RRT*) de Karaman et Frazzoli~\cite{Karaman2011Sampling} construit incrémentalement un arbre de configurations valides et garantit une optimalité asymptotique, c'est-à-dire la convergence vers le chemin optimal à mesure que le nombre d'échantillons croît. Ces planificateurs, dont de nombreuses variantes sont regroupées dans la bibliothèque \textit{Open Motion Planning Library} (OMPL)~\cite{Sucan2012Ompl}, offrent une complétude probabiliste mais produisent des trajectoires géométriquement rugueuses, dont la qualité dépend fortement de la requête et de la scène~\cite{Liu2021Motion} et qui imposent une étape de lissage a posteriori.

\paragraph{Méthodes par optimisation.} À l'inverse, les approches par optimisation locale partent d'une trajectoire initiale et la raffinent en minimisant une fonctionnelle de coût. CHOMP (\textit{Covariant Hamiltonian Optimization for Motion Planning}) de Zucker et al.~\cite{Zucker2013Chomp} exploite le gradient d'un champ de distance signé de la scène, pré-calculé, pour repousser la trajectoire hors des obstacles tout en favorisant sa régularité. TrajOpt de Schulman et al.~\cite{Schulman2014Trajopt} formule plutôt le problème comme une séquence d'optimisations convexes (\textit{sequential convex optimization}) sous contraintes de non-collision : validé sur un jeu de référence de 198 problèmes de planification pour un bras PR2 à sept degrés de liberté, il y atteint des taux de réussite et des temps de résolution supérieurs à ceux de CHOMP et des planificateurs par échantillonnage. Ces méthodes restent toutefois sensibles à la trajectoire d'initialisation et peuvent converger vers des minima locaux sous-optimaux, parfois en violation des contraintes de collision selon la qualité du champ de distance.

\paragraph{Amorçage par modèles génératifs.} Cette dépendance à l'initialisation motive les approches récentes d'amorçage (\textit{warm-starting}), où un modèle appris fournit à l'optimiseur une estimation initiale de bonne qualité : Tian et al.~\cite{Tian2025Warm} et Haffemayer et al.~\cite{Haffemayer2026Warm} montrent qu'un tel amorçage accélère la convergence et améliore la fiabilité des solveurs. Cette idée assure la transition vers les méthodes fondées sur l'apprentissage, qui cherchent à remplacer ou compléter la recherche explicite par des politiques entraînées sur des démonstrations.

\subsection{Apprentissage par démonstration, clonage de comportement et renforcement}
L'apprentissage par démonstration (\textit{Learning from Demonstration}), dont Argall et al.~\cite{Argall2009Survey} dressent la taxonomie fondatrice, vise à reproduire un comportement expert sans programmer explicitement la loi de commande.

\paragraph{Clonage de comportement et erreur cumulative.} Le clonage de comportement (\textit{Behavior Cloning}, BC) traite le problème comme une régression supervisée des actions expertes à partir des observations~\cite{Zhao2023Fine}, une stratégie appliquée avec succès à l'acquisition alimentaire par imitation visuelle~\cite{Sundaresan2022Visuo,Liu2024Adaptive}. Sous une perte quadratique (\textit{Mean Squared Error}, MSE), le BC souffre néanmoins de deux limites structurelles : l'erreur cumulative (\textit{compounding errors}), où de petites déviations éloignent progressivement le robot de la distribution vue à l'entraînement, et l'effondrement de la multimodalité, la moyenne de plusieurs démonstrations valides pouvant produire une action invalide.

\paragraph{Politiques implicites et fondées sur l'énergie.} Pour surmonter ce dernier écueil, Florence et al.~\cite{Florence2021Implicit} proposent le clonage de comportement implicite (\textit{Implicit BC}), qui représente la politique par un modèle d'énergie plutôt que par une régression directe : cette formulation capture des distributions d'actions multivaluées et discontinues et surpasse la régression explicite par MSE sur la majorité des tâches évaluées, jusqu'à des insertions exigeant une tolérance de l'ordre du millimètre. Des travaux connexes formalisent l'imitation strictement hors ligne~\cite{Jarrett2021Strictly} ou exploitent des représentations apprises de façon auto-supervisée pour améliorer la généralisation~\cite{Florence2019Self}.

\paragraph{Apprentissage par renforcement.} L'apprentissage par renforcement (\textit{Reinforcement Learning}, RL) résout nativement la multimodalité par exploration active et permet d'atteindre une manipulation précise~\cite{Luo2025Precise}, y compris sur des tâches à long horizon décomposées en sous-objectifs~\cite{Gupta2019Relay}. Il se heurte cependant à une faible efficacité d'échantillonnage (\textit{sample inefficiency}) et à des risques de collision difficiles à maîtriser lors d'un apprentissage direct sur plateforme physique.

\paragraph{Systèmes dynamiques stables.} Une voie complémentaire encode le mouvement comme un système dynamique dont la stabilité est garantie par construction : l'estimateur de Khansari-Zadeh et Billard~\cite{KhansariZadeh2011Stable} assure ainsi la convergence vers la cible depuis toute condition initiale, au prix d'une expressivité moindre que celle des modèles génératifs récents.

\subsection{Modèles génératifs, politiques de diffusion et Flow Matching}
Les politiques génératives modélisent explicitement la distribution des trajectoires expertes, ce qui leur permet de représenter des comportements multimodaux là où le clonage de comportement classique échoue ; plusieurs revues récentes en dressent le panorama~\cite{WolfXXXXDiffusion,Cho2025Survey}.

\paragraph{Politiques de diffusion.} Héritière des modèles probabilistes de débruitage (\textit{Denoising Diffusion Probabilistic Models}) de Ho et al.~\cite{Ho2020Denoising} et de leur transposition à la planification par Janner et al. avec \textit{Diffuser}~\cite{Janner2022Planning,Huang2023Diffusion}, la \textit{Diffusion Policy} de Chi et al.~\cite{ChiXXXXVisuomotor} génère les actions par débruitage stochastique itératif et améliore en moyenne de 46,9~\% le taux de réussite sur quinze tâches issues de quatre référentiels de manipulation. Sa variante tridimensionnelle, la \textit{3D Diffusion Policy} de Ze et al.~\cite{Ze2024Diffusion}, conditionne la génération sur un nuage de points et surpasse de 24,2~\% les politiques fondées sur l'image avec seulement dix démonstrations. Le principal inconvénient de ces méthodes demeure la latence : le débruitage requiert de nombreuses évaluations du réseau, ce qui limite la fréquence d'inférence en boucle fermée.

\paragraph{Accélération par cohérence et distillation.} Cette latence a suscité une famille de méthodes réduisant le nombre d'étapes d'inférence par distillation ou par contrainte de cohérence~\cite{Prasad2024Consistency,Wang2024One}. ManiCM de Lu et al.~\cite{Lu2025Manicm} impose ainsi une contrainte de cohérence sur le processus de diffusion pour ramener la génération à une seule étape, accélérant l'inférence d'un facteur dix sans perte mesurable du taux de réussite.

\paragraph{Flow Matching conditionnel.} Le \textit{Flow Matching} de Lipman et al.~\cite{Lipman2023Flow} offre une alternative déterministe à la diffusion. L'idée est d'apprendre un champ de vecteurs $v_\theta(\mathbf{x},t)$, indexé par un temps de flux $t\in[0,1]$, qui déplace continûment un échantillon depuis une loi simple $p_0$ (typiquement gaussienne) jusqu'à la distribution des trajectoires expertes $p_1$. Une trajectoire candidate $\mathbf{x}_t$ obéit alors à l'équation différentielle ordinaire (ODE) $\mathrm{d}\mathbf{x}_t/\mathrm{d}t = v_\theta(\mathbf{x}_t, t)$, dont l'intégration de $t=0$ à $t=1$ transforme un bruit initial $\mathbf{x}_0\sim p_0$ en une trajectoire réaliste $\mathbf{x}_1$. La difficulté est que le champ engendrant la distribution marginale reste inconnu ; la contribution du \textit{Conditional Flow Matching}~\cite{Tong2024Conditional} est de montrer qu'on peut l'apprendre en régressant $v_\theta$ sur une cible $u_t(\mathbf{x}\mid\mathbf{x}_1)$ \emph{conditionnée} à un échantillon expert $\mathbf{x}_1$, pour laquelle le déplacement à imiter est connu analytiquement,
\begin{equation}
\mathcal{L}(\theta) = \mathbb{E}_{t,\ \mathbf{x}_1\sim p_1,\ \mathbf{x}\sim p_t(\cdot\mid\mathbf{x}_1)}\Big[\big\| v_\theta(\mathbf{x},t) - u_t(\mathbf{x}\mid\mathbf{x}_1)\big\|^2\Big],
\end{equation}
l'espérance portant sur un temps $t$ tiré uniformément, un échantillon expert $\mathbf{x}_1$ et un point intermédiaire $\mathbf{x}$ du chemin associé. En choisissant un chemin \emph{redressé}~\cite{Liu2023Rectified}, obtenu par interpolation linéaire $\mathbf{x}_t=(1-t)\,\mathbf{x}_0+t\,\mathbf{x}_1$ entre le bruit et la donnée, cette cible se réduit au vecteur constant $\mathbf{x}_1-\mathbf{x}_0$ : le réseau n'a plus qu'à prédire la direction reliant le bruit à la trajectoire experte, ce qui rend les courbes d'intégration quasi rectilignes et permet une génération de haute qualité en très peu d'étapes. Il suffit enfin de conditionner le champ $v_\theta(\mathbf{x},t,\mathbf{o})$ sur une observation $\mathbf{o}$, tel un nuage de points, pour en faire un planificateur génératif ; PointFlowMatch de Chisari et al.~\cite{Chisari2024Manipulation} en a démontré l'application directe à la manipulation à partir de nuages de points.

\paragraph{Variantes accélérées et extensions géométriques.} De nombreuses formulations poussent encore l'efficacité : ManiFlow de Yan et al.~\cite{Yan2025Maniflow} combine \textit{Flow Matching} et entraînement par cohérence pour générer des actions en une à deux étapes tout en doublant presque le taux de réussite réel, tandis que les politiques à vitesse moyenne (\textit{MeanFlow})~\cite{Geng2025Mean,Fang2026Omp} et à flux en continu~\cite{Jiang2025Streaming} visent la génération en une seule étape. Ce cadre a par ailleurs été décliné pour la planification de mouvement de manipulateurs : FlowMP de Nguyen et al.~\cite{Nguyen2025Flowmp} étend le \textit{Flow Matching} aux dynamiques du second ordre pour garantir des trajectoires lisses et physiquement exécutables~\cite{Soleymanzadeh2026Flow}, tandis que la politique riemannienne (\textit{RFMP}) de Braun et al.~\cite{Braun2024Riemannian} respecte la géométrie de l'espace des poses par interpolation géodésique, produisant des trajectoires deux à quatre fois moins saccadées et de 30 à 45~\% plus rapides à l'inférence que les politiques de diffusion. C'est cette combinaison d'expressivité multimodale et d'inférence rapide et déterministe qui motive le choix du \textit{Flow Matching} comme planificateur d'intention nominale dans ce travail.

% =================================================================
\section{Perception visuelle et mémoire spatiale}
% =================================================================

\noindent\textbf{Objectif de la section :} Documenter les mécanismes permettant de convertir un flux de données visuelles bidimensionnelles et bruitées en un signal de conditionnement géométrique stable et persistant.

\subsection{Vision-langage, segmentation zéro-shot et suivi visuel}
La première étape du pipeline consiste à extraire de la scène l'objet pertinent pour la tâche et à l'ancrer sémantiquement. Deux paradigmes s'opposent : les modèles monolithiques, qui fusionnent perception et action, et les architectures modulaires, qui séparent la compréhension sémantique de la génération de mouvement.

\paragraph{Modèles Vision-Langage-Action monolithiques.} Les architectures \textit{Vision-Language-Action} (VLA) fusionnent perception et commande dans un unique réseau entraîné de bout en bout. RT-2 de Brohan et al.~\cite{Brohan2023Vision} et \textit{OpenVLA} de Kim et al.~\cite{Kim2024Openvla} en sont représentatifs : ce dernier, avec sept milliards de paramètres, surpasse le modèle propriétaire RT-2-X de 16,5~\% en taux de réussite absolu sur vingt-neuf tâches tout en mobilisant sept fois moins de paramètres, un domaine dont plusieurs revues retracent l'évolution~\cite{Sapkota2026Vision,Li2025Foundation}. Leur principale limite est une latence d'inférence souvent incompatible avec le contrôle à haute fréquence, que TinyVLA de Wen et al.~\cite{Wen2025Tinyvla} atténue au moyen d'une ossature compacte accélérant l'inférence sans phase de pré-entraînement lourde.

\paragraph{Pipelines modulaires et ancrage sémantique.} À l'opposé, les approches modulaires découplent l'exécution : un modèle de vision-langage interprète l'invite textuelle de la tâche et ancre l'instruction dans la scène~\cite{Liu2024Vision,Mees2022What,Ingelhag2024Skill,Yin2025Codediffuser}, laissant la segmentation fine à un module spécialisé. Ce découplage préserve l'interprétabilité et autorise le remplacement indépendant de chaque composant.

\paragraph{Segmentation promptable et suivi temporel.} Le modèle fondateur \textit{Segment Anything} (SAM) de Kirillov et al.~\cite{Kirillov2023Sam}, entraîné sur le jeu SA-1B de 1,1~milliard de masques répartis sur onze millions d'images, réalise une segmentation zéro-shot souvent comparable aux méthodes pleinement supervisées. Sa troisième version, \textit{Segment Anything 3} (SAM3)~\cite{Meta2025Sam3}, introduit la segmentation par concept (\textit{Promptable Concept Segmentation}) : à partir d'une simple locution nominale, elle retourne d'un coup toutes les instances correspondantes, là où SAM~1 et~2 ne renvoyaient qu'un objet par requête, et double la précision des systèmes antérieurs. Le modèle SAM2 de Ravi et al.~\cite{Ravi2024Sam2} assure quant à lui la propagation temporelle et le suivi du masque sémantique sur le flux de trames, ce qui prépare la reconstruction géométrique de la cible et des obstacles.

\subsection{Reconstruction RGB-D et représentation par nuages de points}
La rétroprojection géométrique des masques sémantiques dans l'espace cartésien, à l'aide d'un capteur RGB-D, transforme les régions image en nuages de points de la cible et des obstacles. Se pose alors la question de la représentation sur laquelle conditionner la politique générative.

\paragraph{Conditionnement sur nuages de points.} Le conditionnement direct sur nuages de points s'est imposé comme une alternative robuste aux observations purement image, car il capture explicitement la structure tridimensionnelle de la scène~\cite{Seita2022Toolflownet,Li2025Language}. La \textit{3D Diffusion Policy} de Ze et al.~\cite{Ze2024Diffusion} l'illustre en surpassant de 24,2~\% les politiques fondées sur l'image sur soixante-douze tâches simulées à partir de dix démonstrations seulement. À plus grande échelle, la politique fondationnelle FP3 de Yang et al.~\cite{Yang2025Fp3}, pré-entraînée sur soixante mille trajectoires couvrant quatorze robots, réussit 42 des 50 tâches inédites évaluées avec un taux de succès de 84~\%, démontrant la généralisation permise par cette représentation.

\paragraph{Robustesse aux occlusions et aux vues uniques.} L'acquisition depuis une caméra unique demeure toutefois sensible aux occlusions induites par l'effecteur. CordViP de Fu et al.~\cite{Fu2025Cordvip} y répond en construisant des nuages de points conscients de l'interaction, à partir d'une estimation de pose 6D et de la proprioception, ce qui rétablit les correspondances objet--main masquées et atteint l'état de l'art sur six tâches réelles ; d'autres travaux suivent plutôt la pose d'objets clés que le nuage brut~\cite{Sun2025Prism}. À l'inverse, certaines études montrent qu'une fusion 2D enrichie de pseudo-3D peut approcher ces performances à moindre coût d'acquisition~\cite{Yu2025Need,Ni2025Semantic}, ce qui nuance l'intérêt systématique de la 3D. Cette représentation capture l'organisation spatiale locale sans requérir de modèles CAO ni de primitives analytiques connus a priori.

\subsection{Cartes persistantes face aux occlusions et pertes de détection}
La perception instantanée ne suffit pas en environnement dynamique : lorsque la cible disparaît temporairement, une couche de mémoire spatiale explicite est nécessaire pour stabiliser les points de conditionnement.

\paragraph{Occlusions en environnement non structuré.} Les occlusions physiques, causées par le corps du robot lui-même ou par les modifications de la scène, perturbent directement l'inférence des politiques visuo-motrices~\cite{Jones2025Beyond}, d'autant plus qu'une politique en boucle ouverte ne dispose d'aucun mécanisme pour combler l'absence transitoire de l'objet.

\paragraph{Cartographie dynamique et estimation de pose incertaine.} La cartographie d'environnements dynamiques offre un cadre établi pour ce problème : DynoSAM de Morris et al.~\cite{Morris2025Dynosam} sépare, dans un même graphe de facteurs, la carte statique du mouvement des objets mobiles, qu'il segmente et reconstruit sans contaminer la structure statique. En complément, SE(3)-PoseFlow de Jin et al.~\cite{Jin2025Poseflow} estime la pose 6D sous forme d'une distribution de probabilité au moyen du \textit{Flow Matching} sur la variété $SE(3)$, capturant la multimodalité inhérente à l'observation partielle là où un réseau déterministe se montrerait exagérément confiant.

\paragraph{Mémoire spatiale persistante.} Sur ces bases, l'intégration de structures cartographiques persistantes (grilles d'occupation historiques, voxelisation locale) permet de conserver la trace spatiale de la cible et des obstacles durant les pertes de détection ou les masquages temporaires, garantissant la continuité du signal de conditionnement fourni au planificateur.

% =================================================================
\section{Représentations géométriques pour la collision}
% =================================================================

\noindent\textbf{Objectif de la section :} Évaluer les méthodes de modélisation de l'encombrement spatial du robot et des obstacles pour formuler des métriques de distance différentiables.

\subsection{Champs de distance de la scène et limites computationnelles}
Imposer une contrainte d'évitement suppose une métrique de distance différentiable entre le robot et les obstacles. Une première famille de méthodes centre cette métrique sur la scène.

\paragraph{Cartographie volumétrique de la scène.} Les champs de distance signés euclidiens (\textit{Euclidean Signed Distance Fields}, ESDF) encodent, en chaque point de l'espace, la distance à l'obstacle le plus proche. Formellement, le champ associe à un point $\mathbf{p}$ sa distance signée à la surface d'un obstacle $\mathcal{O}\subset\mathbb{R}^3$,
\begin{equation}
d(\mathbf{p}) = \operatorname{sgn}(\mathbf{p})\,\min_{\mathbf{q}\in\partial\mathcal{O}}\|\mathbf{p}-\mathbf{q}\|_2,
\end{equation}
comptée négativement à l'intérieur du volume et positivement à l'extérieur, la surface de l'obstacle correspondant au niveau zéro $d(\mathbf{p})=0$. Un tel champ vérifie l'équation eikonale $\|\nabla d\|=1$ : son gradient est unitaire et pointe dans la direction d'éloignement le plus rapide de l'obstacle, ce qui fournit directement la direction d'évitement. Des bibliothèques de cartographie volumétrique telles que OctoMap de Hornung et al.~\cite{Hornung2013Octomap}, Voxblox d'Oleynikova et al.~\cite{Oleynikova2017Voxblox} ou nvblox de Millane et al.~\cite{Millane2024Nvblox} maintiennent un tel champ pour l'évaluation immédiate des collisions, et des transformées de distance calculées sur GPU permettent un évitement réactif à partir d'un nuage de points en direct~\cite{Bedada2024Real}. Leur principale limite est le coût de mise à jour : dès que la scène est dynamique, dense ou fortement bruitée, rafraîchir le champ complet devient trop lent pour alimenter une boucle de correction à haute fréquence.

\subsection{Primitives conservatives et modèles géométriques simplifiés}
Pour contourner ce coût, une seconde famille simplifie plutôt la géométrie du robot lui-même.

\paragraph{Volumes englobants massivement parallèles.} L'approche classique encapsule les liens cinématiques dans des volumes englobants discrets (sphères ou capsules), dont les tests de collision se parallélisent massivement sur GPU. La bibliothèque cuRobo de Sundaralingam et al.~\cite{Sundaralingam2023Curobo} en tire une génération de mouvement en une cinquantaine de millisecondes en moyenne, soit environ soixante fois plus vite que les solveurs d'optimisation de trajectoire de l'état de l'art, et des travaux récents rendent ces enveloppes conscientes de l'outil monté à l'effecteur pour ajuster la marge en temps réel~\cite{Lee2025Fast}. Le revers de cette approximation conservative est qu'elle surestime l'encombrement du robot et réduit l'espace libre exploitable dans les zones confinées.

\subsection{SDF du robot, RDF et approximation par polynômes de Bernstein}
Un changement de paradigme récent déplace le champ de distance de la scène vers le robot, en modélisant directement la géométrie du bras articulé.

\paragraph{Champs de distance du robot (RDF).} Les \textit{Robot Distance Fields} (RDF) de Li et al.~\cite{Li2024rdf} approchent le SDF de chaque lien par des polynômes de Bernstein. Sur une coordonnée spatiale normalisée $\mathbf{u}\in[0,1]^3$, la distance d'un lien $l$ s'écrit comme une combinaison tensorielle de la base de Bernstein de degré $n$,
\begin{equation}
B_{i,n}(u)=\binom{n}{i}\,u^{i}(1-u)^{n-i},\qquad
\bar{d}_l(\mathbf{u})=\sum_{i,j,k} c^{\,l}_{ijk}\,B_{i,n}(u_1)\,B_{j,n}(u_2)\,B_{k,n}(u_3),
\end{equation}
où les polynômes de base $B_{i,n}$ forment une partition lisse de l'unité et où les coefficients $c^{\,l}_{ijk}$, appris hors ligne, encodent la forme du lien $l$. Comme $\bar{d}_l$ est un polynôme, son gradient s'obtient par simple dérivation analytique, sans approximation numérique. Cette formulation fournit ainsi une approximation différentiable du SDF du robot avec une erreur moyenne de l'ordre de 1,4~mm pour un temps de requête de quelques millisecondes et une empreinte mémoire inférieure à 0,03~Mo, là où un SDF neuronal comparable accuse une erreur d'environ 23~mm. Ses gradients étant directement interrogeables par rapport aux points de l'environnement, Govoni et al.~\cite{Govoni2026rdf} les exploitent pour l'évitement réactif de collisions dans un cadre de commande à priorité de tâches.

\paragraph{Alternatives neuronales et limites de régularité.} D'autres approches apprennent plutôt un SDF composite du robot par réseaux de neurones, couplé à de l'optimisation stochastique~\cite{Liu2023Collision}. Dans les deux cas, l'usage d'opérateurs de composition (\textit{min} sur les liens et \textit{soft-min} sur les points) implique que la fonction complète n'est pas strictement $C^\infty$ partout, une propriété à garder à l'esprit lorsque cette distance sert de fonction barrière dans une loi de commande, objet de la section suivante.

% =================================================================
\section{Sécurité réactive et Control Barrier Functions}
% =================================================================

\noindent\textbf{Objectif de la section :} Analyser les lois de commande réactives issues de la théorie du contrôle afin d'imposer une contrainte de sécurité en vitesse sous certaines hypothèses.

\subsection{Champs de potentiel et méthodes réactives classiques}
La sécurité réactive consiste à corriger localement le mouvement à l'exécution. La première formulation historique repose sur des forces virtuelles.

\paragraph{Champs de potentiel artificiels.} Les champs de potentiel artificiels (\textit{Artificial Potential Fields}, APF), introduits par Khatib~\cite{Khatib1986Real}, génèrent une force répulsive à proximité des obstacles et une force attractive vers la cible, dont la résultante dévie le robot en temps réel. Légères en calcul, ces méthodes souffrent toutefois de limites structurelles bien connues : oscillations près des parois, absence de garanties rigoureuses face à des dynamiques complexes et blocage dans les minima locaux des configurations concaves. Le concept reste néanmoins exploité dans des approches génératives récentes, où un champ de potentiel appris guide l'imitation vers des mouvements sûrs~\cite{Ding2025Towards}.

\subsection{Fonctions barrières de contrôle et invariance avant}
Les fonctions barrières de contrôle apportent la rigueur théorique qui manque aux champs de potentiel.

\paragraph{Formulation et garantie d'invariance.} Formalisées par Ames et al.~\cite{Ames2019Control} à partir des certificats de barrière de Wieland et Allgöwer~\cite{Wieland2007Constructive}, les fonctions barrières de contrôle (\textit{Control Barrier Functions}, CBF) définissent l'ensemble sûr comme la région où une fonction $h$ reste positive, $\mathcal{C}=\{\mathbf{q}: h(\mathbf{q})\ge 0\}$, la frontière $h(\mathbf{q})=0$ marquant le contact avec l'obstacle. Cet ensemble est dit invariant vers l'avant si toute trajectoire qui y débute y demeure indéfiniment ; il suffit pour cela que la commande respecte la condition différentielle
\begin{equation}
\nabla h(\mathbf{q})^{\top}\dot{\mathbf{q}} \ge -\kappa\, h(\mathbf{q}), \qquad \kappa>0 .
\end{equation}
L'interprétation est intuitive : loin de l'obstacle ($h$ grand) la contrainte est lâche, mais à mesure que le robot approche la frontière ($h\to 0$) elle force la dérivée $\dot{h}=\nabla h(\mathbf{q})^{\top}\dot{\mathbf{q}}$ à rester au-dessus de zéro, interdisant tout franchissement ; le gain $\kappa$ règle l'agressivité de ce freinage.

\paragraph{Projection de commande par programme quadratique.} Comme cette condition est affine en la vitesse articulaire $\dot{\mathbf{q}}$, elle s'impose commodément à une commande nominale $\dot{\mathbf{q}}_{\mathrm{nom}}$ issue du planificateur au moyen d'un programme quadratique (QP-CBF)~\cite{Ames2017Control} :
\begin{equation}
\dot{\mathbf{q}}^{\star}=\arg\min_{\dot{\mathbf{q}}}\ \tfrac12\big\|\dot{\mathbf{q}}-\dot{\mathbf{q}}_{\mathrm{nom}}\big\|^2
\quad\text{s.c.}\quad \nabla h(\mathbf{q})^{\top}\dot{\mathbf{q}}\ge -\kappa\, h(\mathbf{q}).
\end{equation}
Le filtre cherche donc la vitesse admissible la plus proche de l'intention nominale : tant que celle-ci est sûre, elle passe inchangée ; sinon, elle est corrigée du minimum nécessaire, ce qui revient à une projection sur le demi-espace admissible. Cortez et al.~\cite{Cortez2019Control} ont étendu cette théorie aux systèmes mécaniques de degré relatif deux, où la commande s'exerce au niveau des efforts plutôt que des vitesses.

\subsection{Limites des CBF sous perception bruitée et contrôle discret}
Ces garanties, établies en temps continu et sous modèle parfait, se dégradent lors de l'implémentation physique.

\paragraph{Bruit de perception et échantillonnage discret.} Sur une plateforme réelle, la barrière est évaluée sur des nuages de points bruités et à cadence finie, ce qui motive les formulations en temps discret de la condition barrière proposées par Agrawal et Sreenath~\cite{Agrawal2017Discrete} pour préserver l'invariance malgré la discrétisation.

\paragraph{Faisabilité et actionnement.} S'ajoute le problème de compatibilité des contraintes, lorsque la barrière exige un effort de commande excédant les bornes physiques de l'actionneur~\cite{Beaver2024Optimal}. Enfin, les non-linéarités induites par les approximations géométriques lisses du robot et le suivi assuré par un contrôleur en position bas niveau limitent la validité absolue des preuves d'invariance globale, ce qui invite à traiter la CBF comme un filtre de sécurité pragmatique plutôt que comme une garantie formelle.

% =================================================================
\section{Couplage entre modèles génératifs et filtres de sécurité}
% =================================================================

\noindent\textbf{Objectif de la section :} Examiner les différentes stratégies d'unification de l'IA générative et des lois de commande sécuritaires afin de positionner la méthode de découplage retenue.

Deux grandes stratégies permettent d'unir l'expressivité des modèles génératifs et la rigueur des lois de commande sécuritaires : contraindre la sécurité pendant la génération, ou filtrer la trajectoire au moment de l'exécution.

\subsection{Contraintes intégrées à la génération}
La première stratégie injecte la contrainte au cœur du processus génératif, à chaque étape de l'intégration de l'ODE (guidage intra-ODE).

\paragraph{Barrières et métriques dans le processus génératif.} SafeFlow de Dai et al.~\cite{Dai2025Safeflow} introduit des fonctions barrières de \textit{Flow Matching} (\textit{Flow Matching Barrier Functions}) qui maintiennent la trajectoire dans la région sûre sur tout l'horizon, sans réentraînement. SafeFlowMatcher de Yang et al.~\cite{Yang2026Safeflowmatcher} raffine cette idée par un intégrateur prédiction--correction assorti d'un certificat de barrière garantissant l'invariance avant et la convergence en temps fini. CASF de Long et al.~\cite{Long2026Constraining} convertit quant à lui chaque contrainte en une métrique riemannienne locale, tirée en arrière dans l'espace de commande, qui atténue ou redirige le mouvement au voisinage des frontières tout en préservant la multimodalité des politiques à flux en continu. D'autres variantes imposent la contrainte via des fonctions barrières et de Lyapunov dans le processus de diffusion~\cite{Mizuta2024Cobl,Song2026Omniguide}, via une formulation unifiée admettant égalités et inégalités~\cite{Yang2026Uniconflow}, ou en garantissant l'admissibilité dynamique des trajectoires~\cite{Huang2026Sad}, tandis que certains conditionnent la génération sur la structure de la scène pour amorcer un planificateur sous contraintes~\cite{Haffemayer2026Warm,Mizuta2026Unified}.

\paragraph{Limite : la dérive distributionnelle.} Efficaces en simulation, ces méthodes augmentent le temps de calcul et, surtout, induisent un risque de dérive distributionnelle (\textit{distributional drift}) : en forçant le modèle à traverser des états latents hors de sa distribution d'entraînement, elles peuvent dégrader la qualité et le réalisme de la trajectoire générée.

\subsection{Filtres de sécurité appliqués à l'exécution}
La seconde stratégie laisse la génération intacte et n'intervient qu'en aval, sur la trajectoire finale.

\paragraph{Boucliers cinématiques externes.} Le cadre PACS (\textit{Path-Consistent Safety filtering}) de Römer et al.~\cite{Rmer2026From} en est l'archétype : il applique un freinage cohérent avec la trajectoire générée, vérifié par analyse d'atteignabilité ensembliste, ce qui préserve la distribution apprise tout en garantissant la sécurité en environnement dynamique. Cette approche préserve mieux le succès de la tâche que les filtres réactifs et surpasse les CBF de jusqu'à 68~\% en taux de réussite, en ne modifiant la trajectoire qu'au déploiement, sans perturber la dynamique interne du modèle génératif.

\paragraph{Correction et replanification à l'exécution.} D'autres travaux corrigent localement la politique en cours d'exécution : FlowCorrect de Welte et al.~\cite{Welte2026Flowcorrect} adapte une politique de \textit{Flow Matching} à partir de corrections humaines éparses au moyen d'un module léger, atteignant 80~\% de réussite sur des cas auparavant échoués sans réentraîner le réseau. D'autres encore compensent explicitement le délai d'inférence entre l'observation et l'exécution~\cite{Liao2025Delay}, ou activent une replanification à haute fréquence sans réentraînement~\cite{Ye2025Rapid}.

\subsection{Découplage entre intention nominale et sécurité réactive}
La synthèse de ces deux familles justifie l'architecture retenue dans ce travail. Plutôt que d'entrelacer sécurité et génération, l'analyse précédente motive une isolation fonctionnelle complète entre un planificateur \textit{Flow Matching}, qui porte l'intention globale à long terme sans subir de dérive distributionnelle, et un filtre de sécurité réactif local, agissant en continu à l'exécution~\cite{Rmer2026From,Cho2025Survey}. Ce découplage concilie l'expressivité des distributions apprises et l'impératif de réactivité requis pour intercepter les collisions sur la plateforme réelle, et constitue le fil conducteur de la méthodologie développée au chapitre suivant.

% =================================================================
\section{Synthèse critique et positionnement}
% =================================================================

\noindent\textbf{Objectif de la section :} Résumer les lacunes de l'état de l'art, formaliser l'hypothèse de recherche et présenter un tableau comparatif synthétique des architectures de la littérature.

\subsection{Lacunes identifiées dans la littérature}
L'analyse croisée de la littérature scientifique contemporaine met en exergue trois lacunes majeures :
\begin{enumerate}
    \item \textbf{L'hypothèse d'une perception parfaite :} Les frameworks actuels validant l'alliance du \textit{Flow Matching} et des CBF supposent des obstacles analytiques parfaits connus a priori, se montrant inadaptés face aux nuages de points 3D bruts et bruités.
    \item \textbf{L'absence de persistance temporelle des modèles génératifs :} Les planificateurs génératifs souffrent d'une vulnérabilité cinématique lors d'occlusions temporaires de la cible, faute d'une couche cartographique persistante capable de stabiliser les points de conditionnement.
    \item \textbf{Le coût numérique des représentations globales :} Les approches basées sur la mise à jour de champs de distance globaux ou de \textit{Neural SDF} s'avèrent trop lourdes pour garantir une boucle de rétroaction à haute fréquence face à des obstacles imprévus.
\end{enumerate}

\subsection{Positionnement de l'architecture proposée}
Pour répondre à ces limites, cette thèse soutient la pertinence d'une \textit{architecture découplée entre génération nominale et filtrage réactif de sécurité}. L'intention globale est générée en boucle ouverte par un modèle \textit{Flow Matching} conditionné par une observation géométrique persistante issue de la perception RGB-D et du suivi SAM2, tandis que la sécurité immédiate est déléguée à un filtre de projection CBF articulaire bas niveau s'appuyant sur une approximation polynomiale de Bernstein du SDF local du robot.

\subsection{Hypothèse de recherche}
L'hypothèse centrale de cette recherche s'énonce ainsi :
\begin{quote}
    « L’utilisation d’un filtre SDF-CBF fondé sur un SDF local du robot, approximé par des polynômes de Bernstein et alimenté par une carte perceptuelle persistante, permet de corriger en temps réel les vitesses nominales générées par un modèle Flow Matching, afin d’améliorer l’exécution sécuritaire de trajectoires dans des environnements non structurés et inconnus, sans réentraîner le modèle génératif ni construire un champ de distance global de l’environnement. »
\end{quote}

\vspace{1em}
\noindent Le Tableau~\ref{tab:etat_art} résume le positionnement de l'approche proposée vis-à-vis des frameworks de référence de la littérature.

\begin{table}[h!]
\centering
\caption{Tableau comparatif des frameworks de planification et de sécurité de l'état de l'art.}
\label{tab:etat_art}
\footnotesize
\setlength{\tabcolsep}{2pt}
\renewcommand{\arraystretch}{1.15}
\begin{tabular}{>{\raggedright\arraybackslash}p{2.1cm}>{\raggedright\arraybackslash}p{1.7cm}>{\raggedright\arraybackslash}p{2.1cm}>{\raggedright\arraybackslash}p{2.5cm}>{\raggedright\arraybackslash}p{3.0cm}>{\raggedright\arraybackslash}p{1.5cm}>{\raggedright\arraybackslash}p{2.2cm}}
\hline
\textbf{Cadre / Méthode} & \textbf{Modèle nominal} & \textbf{Sécurité} & \textbf{Représentation} & \textbf{Correction} & \textbf{Occlusions} & \textbf{Réactivité} \\ \hline
Diffusion Policy & Diffusion & Implicite & Espace latent & Aucune / boucle ouverte & Non & Limitée par le débruitage \\
OmniGuide & Flow Matching & Souple & Grille / énergie & Guidage intra-ODE & Non & Moyenne \\
CASF & Streaming Flow & Souple & Neural SDF & Déformation intra-ODE & Partielle & Élevée mais coûteuse \\
SafeFlow\-Matcher & Flow Matching & Explicite CBF & Obstacles analytiques & Correction intra-génération & Non & Élevée en simulation \\
\textbf{Notre approche} & \textbf{Flow Matching} & \textbf{Explicite : projection SDF-CBF} & \textbf{Nuage filtré persistant + SDF robot} & \textbf{Correction vitesse articulaire à l'exécution} & \textbf{Oui, via persistance} & \textbf{Élevée, mesurée expérimentalement} \\ \hline
\end{tabular}
\end{table}

\vspace{1em}
En somme, l'examen critique de la littérature scientifique met en évidence un compromis persistant entre l'expressivité des politiques génératives et la rigueur des algorithmes de préservation de la sécurité. Si le \textit{Flow Matching} lissant et déterministe permet de capturer efficacement la multimodalité des mouvements experts tout en réduisant la latence propre aux processus de diffusion stochastique, son exécution en boucle ouverte au sein d'environnements réels et changeants demeure sujette aux incertitudes perceptives. Les méthodes contemporaines tentant d'incorporer des contraintes au cours de la phase de génération tendent à alourdir le temps de calcul ou à provoquer une dérive distributionnelle nuisible pour le succès de la tâche. Afin d'apporter une solution pragmatique à ces verrous technologiques, le Chapitre 3 détaillera la méthodologie de l'architecture découplée proposée dans ce travail. Nous y exposerons comment un pipeline de perception persistant s'articule de manière synchrone avec un filtre réactif basé sur une approximation par polynômes de Bernstein pour améliorer, en temps réel et sous des contraintes vérifiables localement, l'exécution sécurisée définie au cœur de notre hypothèse de recherche.
